% ══════════════════════════════════════════════════════════════════════
% Supplementary Material
% ══════════════════════════════════════════════════════════════════════
\documentclass[journal]{IEEEtran}

\usepackage[T1]{fontenc}
\usepackage[utf8]{inputenc}
\usepackage{amsmath,amssymb}
\usepackage{graphicx}
\usepackage{booktabs}
\usepackage{multirow}
\usepackage{xcolor}
\usepackage{cite}
\usepackage{threeparttable}
\usepackage{pgfplots}
\pgfplotsset{compat=1.18}
\usepackage{tikz}
\usetikzlibrary{arrows.meta,positioning}
\usepackage[colorlinks=true,citecolor=blue!60!black,linkcolor=blue!60!black]{hyperref}

\begin{document}

\title{Supplementary Material:\\Continual Learning and Constrained RL for Adversarially Robust NIDS}

\author{Roger~Nick~Anaedevha, Alexander~Gennadievich~Trofimov,
        and~Yuri~Vladimirovich~Borodachev}

\maketitle

\section{Dataset Characteristics}

Table~\ref{tab:datasets} provides the detailed characteristics of all six benchmark datasets used in the experimental evaluation.

\begin{table*}[!t]
\centering
\caption{Detailed Dataset Characteristics}
\label{tab:datasets}
\begin{tabular}{@{}llrrrrl@{}}
\toprule
\textbf{Group} & \textbf{Dataset} & \textbf{Flows} & \textbf{Features} & \textbf{Attack Classes} & \textbf{Benign \%} & \textbf{Source} \\
\midrule
\multirow{3}{*}{General Network} & CIC-IoT-2023 & 1{,}621{,}834 & 46 & 33 & 12.4 & CIC, UNB \\
& CSE-CIC-IDS2018 & 16{,}232{,}943 & 79 & 14 & 83.1 & CIC/CSE, AWS \\
& UNSW-NB15 & 2{,}540{,}044 & 49 & 9 & 55.0 & UNSW Canberra \\
\midrule
\multirow{3}{*}{Cloud/Edge} & CloudSec-2024 & 38{,}412{,}201 & 62 & 11 & 71.3 & Kaggle \\
& ContainerIDS & 18{,}764{,}592 & 54 & 8 & 68.7 & Kaggle \\
& EdgeIIoT-2024 & 6{,}628{,}386 & 61 & 8 & 42.1 & Kaggle \\
\midrule
& \textbf{Total} & \textbf{84{,}200{,}000} & --- & \textbf{83 unique} & --- & --- \\
\bottomrule
\end{tabular}
\end{table*}

\section{Sequential Task Split Protocol}

For each dataset, the five temporal splits are constructed as follows. Flows are first sorted by timestamp (or by index order for datasets without explicit timestamps). The sorted corpus is then divided into five equal-sized partitions, each representing a simulated deployment period. Within each partition, an 80/20 stratified train/test split is applied, preserving the class distribution of that period. The held-out test set for evaluation metrics spans all five periods uniformly.

For cross-dataset sequential evaluation, the six datasets are presented in the following order: CIC-IoT-2023 (33~classes), UNSW-NB15 (9~classes, 7~novel), CSE-CIC-IDS2018 (14~classes, 5~novel), CloudSec-2024 (11~classes, 4~novel), ContainerIDS (8~classes, 2~novel), and EdgeIIoT-2024 (8~classes, 2~novel). The ``novel'' count indicates classes not encountered in any prior dataset in the sequence.

\section{Per-Class Accuracy Retention Matrix}

Table~\ref{tab:perclass_full} reports the per-class accuracy retention on CIC-IoT-2023 after training through all five sequential splits, comparing fine-tuning against the proposed EWC\,+\,Replay method. Accuracy retention is defined as $a_{5,k} / a_{1,k}$ for class $k$.

\begin{table}[!t]
\centering
\caption{Per-Class Accuracy Retention on CIC-IoT-2023 (\%)}
\label{tab:perclass_full}
\begin{tabular}{@{}lcc@{}}
\toprule
\textbf{Attack Family} & \textbf{Fine-Tuning} & \textbf{EWC+Replay} \\
\midrule
DDoS-UDP Flood & 82.6 & 98.3 \\
DDoS-ICMP Flood & 84.1 & 98.7 \\
DDoS-SYN Flood & 80.3 & 97.9 \\
DDoS-TCP Flood & 79.8 & 97.6 \\
DDoS-PSHACK Flood & 77.2 & 97.1 \\
DDoS-HTTP Flood & 73.4 & 96.2 \\
DoS-UDP Flood & 81.5 & 97.8 \\
DoS-SYN Flood & 80.1 & 97.4 \\
DoS-TCP Flood & 78.9 & 97.2 \\
DoS-HTTP Flood & 74.6 & 96.5 \\
Recon-PortScan & 82.3 & 96.8 \\
Recon-OSScan & 79.4 & 96.1 \\
Recon-HostDiscovery & 77.8 & 95.4 \\
Recon-PingSweep & 76.2 & 95.0 \\
Mirai-UDP & 83.7 & 98.1 \\
Mirai-ACK & 82.1 & 97.6 \\
Mirai-SYN & 81.4 & 97.3 \\
Mirai-HTTP & 78.6 & 96.8 \\
Mirai-Greeth & 76.9 & 96.2 \\
Mirai-Greip & 75.3 & 95.8 \\
BruteForce & 71.2 & 93.5 \\
BF-SSH & 69.8 & 92.8 \\
Spoofing-ARP & 68.4 & 92.1 \\
Spoofing-DNS & 66.7 & 91.4 \\
Web-XSS & 70.1 & 93.2 \\
Web-SQLi & 68.4 & 92.8 \\
Web-CommandInj & 67.2 & 91.9 \\
Web-Backdoor & 65.8 & 91.1 \\
VulnScan & 73.6 & 94.6 \\
DictionaryBF & 72.1 & 93.8 \\
Uploading & 69.5 & 92.4 \\
MITM & 71.8 & 93.1 \\
Fingerprinting & 74.2 & 94.2 \\
Benign & 82.6 & 98.4 \\
\midrule
\textbf{Average} & \textbf{76.1} & \textbf{95.7} \\
\bottomrule
\end{tabular}
\end{table}

\section{Drift Detection Evaluation}

The KL-divergence drift detector is evaluated against two types of drift events: synthetic drift (created by resampling class distributions to simulate sudden prevalence changes) and natural drift (observed between temporal splits within CIC-IoT-2023). Table~\ref{tab:drift_eval} reports the detection performance.

\begin{table}[!t]
\centering
\caption{Drift Detection Performance}
\label{tab:drift_eval}
\begin{tabular}{@{}lcc@{}}
\toprule
\textbf{Metric} & \textbf{Synthetic Drift} & \textbf{Natural Drift} \\
\midrule
True Positive Rate & 0.96 & 0.91 \\
False Positive Rate & 0.04 & 0.08 \\
AUC & 0.97 & 0.92 \\
Mean Detection Delay (batches) & 1.2 & 2.4 \\
\bottomrule
\end{tabular}
\end{table}

The detector achieves higher accuracy on synthetic drift because the class-marginal shifts are more pronounced. Under natural drift, where changes may be gradual and multi-modal, the mean detection delay increases to 2.4~batches, reflecting the inherent difficulty of detecting slow concept drift via marginal KL-divergence.

\section{Hyperparameter Sensitivity: Full Results}

\subsection{EWC Consolidation Strength $\lambda$}

Table~\ref{tab:lambda_sensitivity} reports the sensitivity of average accuracy and backward transfer to $\lambda$ across all six datasets.

\begin{table}[!t]
\centering
\caption{Sensitivity to $\lambda$ (Average Across Datasets)}
\label{tab:lambda_sensitivity}
\begin{tabular}{@{}rcc@{}}
\toprule
$\lambda$ & \textbf{AA (\%)} & \textbf{BWT} \\
\midrule
100 & 92.1 & $-$0.051 \\
500 & 94.3 & $-$0.033 \\
1{,}000 & 95.6 & $-$0.022 \\
\textbf{5{,}000} & \textbf{96.8} & \textbf{$-$0.014} \\
10{,}000 & 96.4 & $-$0.016 \\
15{,}000 & 95.2 & $-$0.024 \\
\bottomrule
\end{tabular}
\end{table}

\subsection{Replay Buffer Size $M$}

Table~\ref{tab:buffer_sensitivity} reports the sensitivity to replay buffer size.

\begin{table}[!t]
\centering
\caption{Sensitivity to Buffer Size $M$ (CIC-IoT-2023)}
\label{tab:buffer_sensitivity}
\begin{tabular}{@{}rccc@{}}
\toprule
$M$ & \textbf{AA (\%)} & \textbf{BWT} & \textbf{Memory (MB)} \\
\midrule
500 & 94.2 & $-$0.031 & 2.4 \\
1{,}000 & 95.1 & $-$0.024 & 4.8 \\
3{,}000 & 96.5 & $-$0.016 & 14.4 \\
\textbf{5{,}000} & \textbf{97.1} & \textbf{$-$0.012} & \textbf{24.0} \\
10{,}000 & 97.4 & $-$0.010 & 48.0 \\
\bottomrule
\end{tabular}
\end{table}

\section{RL Training Curves}

The CPO agent's training dynamics are characterised by three key quantities: cumulative reward, false-positive cost, and constraint satisfaction. During the first $5\times10^5$ steps, the agent explores aggressively, with the false-positive cost temporarily exceeding the threshold $\epsilon_{\mathrm{fp}}$. The CPO constraint enforcement mechanism corrects this by increasing the Lagrange multiplier, and by step $10^6$ the constraint is satisfied consistently. By contrast, unconstrained PPO never reliably satisfies the false-positive constraint, oscillating around a violation rate of 2--3\% throughout training.

The Lagrangian PPO baseline shows characteristic oscillation: the dual variable (Lagrange multiplier) increases when the constraint is violated, suppressing aggressive actions, but then decreases as the policy becomes too conservative, permitting a new cycle of violations. CPO's trust-region approach avoids this oscillation by providing monotonic constraint improvement guarantees within each update step.

\section{RL Action Distribution Analysis}

Table~\ref{tab:action_dist} reports the fraction of actions selected by each policy across the evaluation episodes.

\begin{table}[!t]
\centering
\caption{Action Distribution Across Policies (\%)}
\label{tab:action_dist}
\begin{tabular}{@{}lcccc@{}}
\toprule
\textbf{Action} & \textbf{Rule} & \textbf{PPO} & \textbf{L-PPO} & \textbf{CPO} \\
\midrule
Monitor & 62.1 & 38.4 & 44.2 & 41.8 \\
RateLimit & 18.3 & 26.7 & 28.4 & 30.1 \\
Reset & 10.2 & 18.6 & 16.2 & 17.4 \\
Block & 7.1 & 12.8 & 8.9 & 8.6 \\
Quarantine & 2.3 & 3.5 & 2.3 & 2.1 \\
\bottomrule
\end{tabular}
\end{table}

The CPO agent uses \textsc{RateLimit} more frequently than other RL baselines (30.1\% versus 26.7\% for PPO), reflecting a learned preference for moderate-severity actions that provide meaningful mitigation without the high false-positive risk of \textsc{Block} and \textsc{Quarantine}. Unconstrained PPO uses \textsc{Block} 48\% more frequently than CPO (12.8\% versus 8.6\%), which directly accounts for its elevated false-positive rate.

\section{Cross-Dataset Sequential Results}

Table~\ref{tab:crossdataset} reports the continual learning metrics when the six datasets are presented sequentially, simulating a deployment encountering fundamentally different network environments over time.

\begin{table}[!t]
\centering
\caption{Cross-Dataset Sequential Evaluation}
\label{tab:crossdataset}
\begin{tabular}{@{}lccc@{}}
\toprule
\textbf{Method} & \textbf{AA (\%)} & \textbf{BWT} & \textbf{FWT} \\
\midrule
Full Retrain & 95.6 & 0.000 & --- \\
Fine-Tuning & 74.2 & $-$0.218 & +0.061 \\
EWC Only & 87.3 & $-$0.074 & +0.034 \\
GEM & 89.8 & $-$0.051 & +0.021 \\
\textbf{EWC+Replay} & \textbf{93.1} & \textbf{$-$0.024} & \textbf{+0.041} \\
\bottomrule
\end{tabular}
\end{table}

The cross-dataset scenario is substantially more challenging than within-dataset sequential learning because each new dataset introduces entirely different feature distributions, class taxonomies, and traffic patterns. Fine-tuning exhibits catastrophic forgetting with BWT of $-0.218$, indicating that performance on the first dataset (CIC-IoT-2023) degrades by over 21~percentage points after training on all subsequent datasets. The proposed method limits this degradation to 2.4~points, maintaining 93.1\% average accuracy, within 2.5~points of the full-retraining upper bound.

\section{Statistical Significance}

All reported metrics are averaged over five independent runs with different random seeds. Statistical significance is assessed via paired $t$-tests comparing the proposed EWC\,+\,Replay method against each baseline. All reported improvements are significant at $p < 0.01$ for the primary metrics (AA and BWT), with the exception of the unified FIM ablation (Section~V-E of the main paper), where the improvement over independent FIM computation is significant only at $p < 0.05$.

\end{document}
